\sourcepage{162}
\begin{references}
\item Lehr R. Sixteen S-squared over D-squared: a relation for crude sample size estimates. Stat Med. 1992;11:1099-1102.
\item Li H, Wang L, Wei L, Quan H. Sample size calculation for count data in comparative clinical trials with nonuniform patient accrual and early dropout. J Biopharm Stat. 2015;25(1):1-15.
\item Barthel FM, Babiker A, Royston P, Parmar MK. Evaluation of sample size and power for multi-arm survival trials allowing for non-uniform accrual, non- proportional hazards, loss to follow-up and cross-over. Stat Med. 2006;25(15):2521-2542.
\item Ahnn S, Anderson SJ. Sample size determination in complex clinical trials comparing more than two groups for survival endpoints. Stat Med. 1998;17(21):2525-2534.
\item Kerry SM, Bland JM. Trials which randomize practices II: sample size. Fam Pract. 1998;15:84-87.
\item Hemming K, Girling AJ, Sitch AJ, et al. Sample size calculations for cluster randomised controlled trials with a fixed number of clusters. BMC Med Res Methodol. 2011;11:102.
\item Jahn-Eimermacher A, Ingel K, Schneider A. Sample size in cluster-randomized trials with time to event as the primary endpoint. Stat Med. 2013;32(5):739-751.
\item Edwardes MD. Sample size requirements for case-control study designs. BMC Med Res Methodol. 2001;1:11.
\item Drescher K, Timm J, Jöckel KH. The design of case-control studies: the effect of confounding on sample size requirements. Stat Med. 1990;9:765-776.
\item Lui KJ. Sample size determination for case-control studies: the influence of the joint distribution of exposure and confounder. Stat Med. 1990;9:1485-1493.
\item Latouche A, Porcher R, Chevret S. Sample size formula for proportional hazards modelling of competing risks. Stat Med. 2004;23(21):3263-3274.
\item Novikov I, Fund N, Freedman LS. A modified approach to estimating sample size for simple logistic regression with one continuous covariate. Stat Med. 2010;29(1):97-107.
\item Vaeth M, Skovlund E. A simple approach to power and sample size calculations in logistic regression and Cox regression models. Stat Med. 2004;23(11):1781-1792.
\item Dupont WD, Plummer WD Jr. Power and sample size calculations for studies involving linear regression. Control Clin Trials. 1998;19:589-601.
\sourcepage{163}
\item Murcray CE, Lewinger JP, Conti DV, et al. Sample size requirements to detect gene-environment interactions in genome-wide association studies. Genet Epidemiol. 2011;35(3):201-210.
\item Wang S, Zhao H. Sample size needed to detect gene-gene interactions using linkage analysis. Ann Hum Genet. 2007;71(Pt 6):828-842.
\item Witte JS. Rare genetic variants and treatment response: sample size and analysis issues. Stat Med. 2012;31(25):3041-3050.
\item Willan AR. Sample size determination for cost-effectiveness trials. Pharmacoeconomics. 2011;29(11):933-949.
\item Glick HA. Sample size and power for cost-effectiveness analysis (Part 2): the effect of maximum willingness to pay. Pharmacoeconomics. 2011;29(4):287-296.
\item Glick HA. Sample size and power for cost-effectiveness analysis (Part 1). Pharmacoeconomics. 2011;29(3):189-198.
\item Patel HI. Sample size for a dose-response study. J Biopharm Stat. 1992;2:l- 8.
\item Day SJ, Graham DF. Sample size estimation for comparing two or more treatment groups in clinical trials. Stat Med. 1991;10:33-43.
\item Guo JH, Chen HJ, Luh WM. Sample size planning with the cost constraint for testing superiority and equivalence of two independent groups. Br J Math Stat Psychol. 2011;64(3):439-461.
\item Zhang P. A simple formula for sample size calculation in equivalence studies. J Biopharm Stat. 2003;13(3):529-538.
\item Stucke K, Kieser M. A general approach for sample size calculation for the three-arm ‘gold standard’ non-inferiority design. Stat Med. 2012;31(28):3579-3596.
\item Julious SA, Owen RJ. A comparison of methods for sample size estimation for non-inferiority studies with binary outcomes. Stat Methods Med Res. 2011;20(6):595-612.
\item Obuchowski NA. Sample size tables for receiver operating characteristic studies. AJR Am J Roentgenol. 2000;175(3):603-608.
\item Simel DL, Samsa GP, Matchar DB. Likelihood ratios with confidence: sample size estimation for diagnostic test studies. J Clin Epidemiol. 1991;44:763-770.
\item Sim J, Wright CC. The kappa statistic in reliability studies: use, interpretation, and sample size requirements. Phys Ther. 2005;85(3):257-268.
\item Jewell NP. Statistics for Epidemiology. Chapman and Hall; 2004:68.
\end{references}
